\clearpage\pagestyle{otMain}\setlength{\otBlocEspacement}{2.5mm}\fontsize{11}{15}\selectfont\corrigesChapitre{4}{Additionner et soustraire des relatifs}
\begin{blocLivre}{cor-4-1}\begin{corrige}{1}{Même signe positif}\demarche{Additionner les distances à zéro.}
Les termes sont positifs : $3+8=\res{11}$.
\end{corrige}\end{blocLivre}
\begin{blocLivre}{cor-4-2}\begin{corrige}{2}{Même signe négatif}\demarche{Conserver le signe commun.}
On additionne 4 et 9 puis on garde le signe négatif : $\res{-13}$.
\end{corrige}\end{blocLivre}
\begin{blocLivre}{cor-4-3}\begin{corrige}{3}{Signes contraires}\demarche{Comparer les distances à zéro.}
12 est la plus grande distance ; son signe est négatif. $12-7=5$, donc $\res{-5}$.
\end{corrige}\end{blocLivre}
\begin{blocLivre}{cor-4-4}\begin{corrige}{4}{Soustraire un négatif}\demarche{Ajouter l’opposé.}
Soustraire $-8$ revient à ajouter son opposé 8 : $6+8=\res{14}$.
\end{corrige}\end{blocLivre}
\begin{blocLivre}{cor-4-5}\begin{corrige}{5}{Deux opposés}\demarche{Repérer les opposés.}
Les deux nombres sont opposés : leur somme vaut $\res0$.
\end{corrige}\end{blocLivre}
\begin{blocLivre}{cor-4-6}\begin{corrige}{6}{Deux décimaux}\demarche{Séparer signe et calcul décimal.}
$4{,}15-2{,}80=\res{1{,}35}$. Le second calcul vaut $-1{,}6+(-2{,}9)=\res{-4{,}5}$.
\end{corrige}\end{blocLivre}
\begin{blocLivre}{cor-4-7}\begin{corrige}{7}{Réorganiser une somme}\demarche{Chercher les paires d’opposés.}
On regroupe $-17+17=0$ et $6-6=0$. Il reste $\res{-4}$.
\end{corrige}\end{blocLivre}
\begin{blocLivre}{cor-4-8}\begin{corrige}{8}{Parenthèses utiles}\demarche{Identifier l’objet de la soustraction.}
Dans le premier : $-2+7=5$, donc $5-5=\res0$. Dans le second : $5+2+7=\res{14}$. Les parenthèses changent ce qu’on soustrait.
\end{corrige}\end{blocLivre}
\begin{blocLivre}{cor-4-9}\begin{corrige}{9}{Écrire plus simplement}\demarche{Transformer chaque soustraction.}
On obtient $-8+3+6-4$. Les positifs totalisent 9, les négatifs $-12$ ; le résultat vaut $\res{-3}$.
\end{corrige}\end{blocLivre}
\begin{blocLivre}{cor-4-10}\begin{corrige}{10}{Une erreur de signe}\demarche{Contrôler par l’opération inverse.}
Il faut $-5+9=\res4$. Vérification par l’addition inverse : $4+(-9)=-5$. Le résultat $-14$ correspondrait à $-5+(-9)$.
\end{corrige}\end{blocLivre}
\begin{blocLivre}{cor-4-11}\begin{corrige}{11}{Températures}\demarche{Suivre les variations dans l’ordre.}
On additionne les variations : $-4{,}5+7-3{,}2=\res{-0{,}7\,^\circ\mathrm C}$. Le passage par $2{,}5$ degrés permet de contrôler chaque étape.
\end{corrige}\end{blocLivre}
\begin{blocLivre}{cor-4-12}\begin{corrige}{12}{Un ascenseur}\demarche{Distinguer position et déplacement.}
Il arrive à $-3+8-6+2=\res1$. Sa variation totale est $8-6+2=\res4$ niveaux vers le haut. Contrôle : $-3+4=1$.
\end{corrige}\end{blocLivre}
\begin{blocLivre}{cor-4-13}\begin{corrige}{13}{Distance sur une droite}\demarche{Faire valeur de droite moins valeur de gauche.}
B est à droite de A. La distance vaut $2{,}25-(-6{,}5)=2{,}25+6{,}5=\res{8{,}75}$ unités. Elle est positive.
\label{fig-cor-4-13}\figureGeo[corrige]{add-13}
\end{corrige}\end{blocLivre}
\begin{blocLivre}{cor-4-14}\begin{corrige}{14}{Retrouver une variation}\demarche{Soustraire la valeur initiale à la finale.}
La variation vaut $7{,}25-(-18{,}40)=\res{25{,}65}$ euros. Il s’agit d’un apport de cette somme, car la variation est positive.
\end{corrige}\end{blocLivre}
\begin{blocLivre}{cor-4-15}\begin{corrige}{15}{Un bilan à reconstituer}\demarche{Remonter les opérations en sens inverse.}
On remonte les actions : avant la dernière perte, le score était $-6+9=3$ ; avant le gain, $3-7=-4$ ; au départ, $-4+12=\res8$. Contrôle : $8-12+7-9=-6$.
\end{corrige}\end{blocLivre}
